\documentclass[12pt]{article}
\usepackage{resource/shortex}
\input{resource/teaching}

%--------------------
\begin{document}

\begin{center}
\textbf{\Large Project 2 \\  Statistical Data Analysis } 
\end{center}


\section*{Overview}
Your group will select dataset(s) of your interest and conduct data analyses.
This should include :
\begin{itemize}
    \item One numerical data anlysis
    \item One categorical data analysis 
\end{itemize} 
You will write a report of your analyses and submit it together with your R code.
The report should be approximately 5-10 pages in length (12-point font, single-spaced), including references, R code and data file, but the page count is flexible.


\section*{Suggested outline for your writeup}
\begin{itemize}
    \item Introduction
    \begin{enumerate}
        \item Introduce the dataset and its source
        \item Describe the context and background of the data
        \item  Define your variables of interest
        \item Tell the story of your reasoning and possible hypothesis
    \end{enumerate}
    \item Categorical Data Analysis Part
    \begin{enumerate}
        \item Preliminary Data Exploration (EDA) 
            \begin{enumerate}
            \item Summarize categorical variables
            \item Contingency Table 
            \item Visualizations: bar plots, pie charts, etc.
            \end{enumerate} 
        \item Inference for Categorical Data
        \begin{enumerate}
            \item Conditions of inference
            \item Hypoetheses
            \item Test statistic and p-value
            \item Conclusion : Discuss results in context of the data
        \end{enumerate}
    \end{enumerate}
    \item Numerical Data Analysis Part
        \begin{enumerate}
             \item Preliminary Data Exploration (EDA) 
            \begin{enumerate}
                \item Summarize numerical variables (mean, median, sd, IQR, etc)
                \item Visualizations: histograms, boxplots, etc.
            \end{enumerate}
            \item Inference for Numerical Data
        \begin{enumerate}
            \item Conditions of inference
            \item Hypoetheses
            \item Test statistic and p-value
            \item Conclusion : Discuss results in context of the data
        \end{enumerate}
    \end{enumerate}
    \item Conclusion
    \begin{enumerate}
            \item Summarize key insights
            \item Suggest future analysis or potential applications
        \end{enumerate} 
\end{itemize}

\section*{Project requirements}
\begin{itemize}
    \item Source of data must be cited
    \item R code should be included (Rscript (.R) or Rmd files)
    \item References should be provided at the end
    \item Using certain functions in R are not allowed
    \begin{itemize}
        \item Using functions that directly perform hypothesis tests (e.g., t.test(), chisq.test(), prop.test(), etc.)
        \item If you want to use ANOVA, you can use aov(), anova() or lm() functions.
    \end{itemize}
\end{itemize}



\section*{Possible Data Sources to Consider (not limited to these)}
You can click the links below to explore possible data sources.
\begin{itemize}
    \item \href{https://www.openintro.org/data/}{OpenIntro Datasets}
    \item \href{https://cran.r-project.org/web/packages/fivethirtyeight/vignettes/fivethirtyeight.html}{fivethirtyeight 
    R Package}
    \item \href{https://stat.ethz.ch/R-manual/R-devel/library/datasets/html/00Index.html}{datsets R package}
    \item \href{https://www.kaggle.com/datasets}{Kaggle datasets}
    \item \href{https://data.gov/}{US government's open data}
\end{itemize}


\section*{Some example of analyses you can do}
\begin{itemize}
     \item Categorical Data Analysis Examples (choose at least one)
    \begin{itemize}
        \item One categorical data with binary outcome (Inference for a single proportion : Chapter 6.1)
        \item Two Categorical variables with binary outcomes (Inference for two proportions : Chapter 6.2)
        \item One categorical variable with more than two outcomes (Chi-square goodness-of-fit test : Chapter 6.3)
        \item Two categorical variables with more than two outcomes (Testing for independece in two-way tables : Chapter 6.4)
    \end{itemize}

    \item Numerical Data Analysis Examples (choose at least one)
    \begin{itemize}
        \item One numerical variable (One-sample means with t-test : Chapter 7.1)
        \item Two numerical variables (Paired data or Two-sample means with t-test : Chapters 7.2, 7.3)
        \item More than two numerical variables (Comparing many means with ANOVA : Chapter 7.5)
    \end{itemize}
   
\end{itemize}

\section*{Project Timeline}
\begin{enumerate}
    \item Project 2 logistic week : Week 8 (on March 18th)
    \item Deliverable 1 : Project Outline (by March 24th) 
    \item Project 2 work 2 (on April 1st)
    \item Deliverable 2 : Project Plan (Results slides to explain + R code + Datafile) (by April 7th)
    \item Project 2 work 3 (on April 9th)
    \item Deliverable 3 : Writeup + R code(final version) + Datafile (by April 15th)
\end{enumerate}

\section*{Example Roles by Tasks}
\begin{enumerate}
    \item For a group of 4 \\
    Role 1 : Introduction + Categorical Data Analysis   \\
    Role 2 : Categorical Data Analysis \\
    Role 3 : Numerical Data Analysis \\
    Role 4 : Conclusion + Numerical Data Analysis  \\
    \vspace{1em}

    \item For a group of 3 \\
    Role 1 : Introduction, EDA part for numerical and categorical analysis \\
    Role 2 : Numerical Data Analysis \\
    Role 3 : Categorical Data Analysis + Conclusion \\
\end{enumerate}

\section*{Deliverables}
\begin{itemize}
    \item Deliverable 1 : Project Outline (by March 24th) 
    \item Deliverable 2 : Project Plan (Results slides + R code + Datafile) (by April 7th) 
    \item Deliverable 3 : Writeup + R code(final version + Datafile) (by April 15th)
\end{itemize}



\end{document}



% shall i add abstract part too? 
